\documentclass[../../main.tex]{subfiles}
\begin{document}

\begin{center}
  {\fontsize{48pt}{58pt}\selectfont \textbf{1987}\normalsize\textbf{年度}}
\end{center}

\begin{table}[htb]
  \centering
  \caption{出題テーマと難易度}
  \begin{tabular}{|c|p{9cm}|c|}
    \hline
    \textbf{問題番号} & \textbf{テーマ} & \textbf{難易度} \\
    \hline
    第1問 & 多項式が同じ二次式で割り切れる条件 & \\
    \hline
    第2問 & 行列の漸化式で定まる数列 & \\
    \hline
    第3問 & 三角関数の不等式を満たす係数の領域 & \\
    \hline
    第4問 & 三次曲線の接線の本数と弦の長さ & \\
    \hline
    第5問 & 平面による球の切断と球欠の体積 & \\
    \hline
    第6問 & 札の交換で元の状態に戻る確率 & \\
    \hline
  \end{tabular}
\end{table}

\begin{table}[htb]
  \centering
  \caption{解答}
  \renewcommand{\arraystretch}{1.4}
  \begin{tabular}{|c|c|p{9cm}|}
    \hline
    \textbf{問題番号} & \textbf{小問} & \textbf{答え} \\
    \hline
    第1問 & -- & $f,g$がともに$h$で割り切れるのは$a=1,c=0$のとき．それ以外はともに割り切れない \\
    \hline
    第2問 & -- &  \\
    \hline
    第3問 & -- & $(a,b)$の領域は解答の図を参照（境界を含まない） \\
    \hline
    第4問 & (1) & $a^2-3(b-m)$が正なら$2$本，零なら$1$本，負なら$0$本 \\
    \cline{2-3}
    & (2) & $P_1Q_1=P_2Q_2$（証明問題） \\
    \hline
    第5問 & (1) & $\dfrac{(2-a)(a,1,1)}{a^2+2}$ \\
    \cline{2-3}
    & (2) & $a=-1$で最小体積$\pi(18-8\sqrt{3})$ \\
    \hline
    第6問 & (1) & 起こりうる状態は解答を参照 \\
    \cline{2-3}
    & (2) & 遷移確率は解答を参照 \\
    \cline{2-3}
    & (3) & $q_n=\dfrac{2}{5}+\dfrac{3}{5}(-\dfrac{1}{4})^n$ \\
    \hline
  \end{tabular}
\end{table}

\paragraph{試験形態}
6問．試験時間は要確認．

\end{document}
